% Part: lambda-calculus
% Chapter: syntax
% Section: de-bruijn

\documentclass[../../../include/open-logic-section]{subfiles}

\begin{document}

\olfileid{lam}{syn}{deb}

\olsection{દે બ્રુઇન સૂચકાંક}

$\alpha$-સામ્ય ખૂબ સાહજિક છે, કારણ કે $\alpha$-સામ્યવાળાં
પદોનો ``અર્થ એક જ'' હોય છે. હકીકતમાં, લૅમ્બડા પદો માટે એવી
વાક્યરચના આપી શકાય જેમાં $\alpha$-રૂપાંતરણથી એકબીજામાં
બદલી શકાય તેવાં પદો વચ્ચે ભેદ ન રહે. એમાં સૌથી જાણીતી
રીત ચલોની જગ્યાએ તેમના \emph{દે બ્રુઇન સૂચકાંક} વાપરે છે.

$\lambd[x][M]$ લખીએ ત્યારે વિધેયનો પ્રાચલ~$x$ છે એમ
સ્પષ્ટ કહીએ છીએ, જેથી~$M$માં~$x$ વડે તે પ્રાચલનો નિર્દેશ
કરી શકીએ. પરંતુ દે બ્રુઇન સૂચકાંક પદ્ધતિમાં પ્રાચલોનાં નામ
નથી. વિધેયના મુખ્ય ભાગમાં તેમનો નિર્દેશ, પ્રાચલની ઘટના અને
તેને બાંધતા અમૂર્તીકરણ વચ્ચે અમૂર્તીકરણનાં કેટલાં સ્તર છે
તે દર્શાવતી સંખ્યા વડે થાય છે. દાખલા તરીકે,
$\lambd[x][\lambd[y][y x]]$ જુઓ: બહારનું અમૂર્તીકરણ
ચલ~$x$ને બાંધે છે; અંદરનું અમૂર્તીકરણ ચલ~$y$ને બાંધે છે.
ઉપપદ~$y x$ અંદરના અમૂર્તીકરણના વ્યાપમાં છે. $y$ અને
તેને બાંધતા~$\lambd[y]$ વચ્ચે કોઈ અમૂર્તીકરણ નથી, પરંતુ
$x$ અને તેને બાંધતા~$\lambd[x]$ વચ્ચે એક અમૂર્તીકરણ છે.
આથી~$y x$ માટે~$0\, 1$ અને આખા પદ માટે
$\lambd[][\lambd[][01]]$ લખીએ છીએ.

\begin{defn}
  દે બ્રુઇન પદો આગમનથી આ રીતે વ્યાખ્યાયિત થાય છે:
  \begin{enumerate}
  \item $n$, જ્યાં~$n$ કોઈપણ પ્રાકૃતિક સંખ્યા છે.
  \item $PQ$, જ્યાં~$P$ અને~$Q$ બંને દે બ્રુઇન પદો છે.
  \item $\lambd[][N]$, જ્યાં~$N$ દે બ્રુઇન પદ છે.
  \end{enumerate}
\end{defn}

સામાન્ય લૅમ્બડા પદોમાંથી દે બ્રુઇન સૂચકાંકવાળાં પદોમાં
ઔપચારિક રૂપાંતરણ આ પ્રમાણે છે:
\begin{defn}
  \begin{align*}
    F_\Gamma(x) &= \Gamma(x) \\
    F_\Gamma(PQ) &= F_\Gamma(P)F_\Gamma(Q) \\
    F_\Gamma(\lambd[x][N]) &= \lambd[][F_{x,\Gamma}(N)]
  \end{align*}
  અહીં~$\Gamma$ એ શૂન્યથી સૂચકાંકિત કરેલી ચલોની યાદી છે,
  અને~$\Gamma(x)$ એ~$\Gamma$માં ચલ~$x$નું સ્થાન દર્શાવે છે.
  દાખલા તરીકે, $\Gamma$ એ~$x,y,z$ હોય તો~$\Gamma(x)$ એ
  $0$ અને~$\Gamma(z)$ એ~$2$ છે.

  $x,\Gamma$ એટલે~$\Gamma$ની શરૂઆતમાં~$x$ ઉમેરીને મળતી
  યાદી. અગાઉના ઉદાહરણમાં~$\Gamma$ માટે
  $w,\Gamma$ એ~$w,x,y,z$ છે.
\end{defn}

દે બ્રુઇન પદમાંથી સામાન્ય લૅમ્બડા પદ આ રીતે ફરી મેળવી શકાય:

\begin{defn}
  \begin{align*}
    G_\Gamma(n) &= \Gamma[n] \\
    G_\Gamma(PQ) &= G_\Gamma(P) G_\Gamma(Q) \\
    G_\Gamma(\lambd[][N]) &= \lambd[x][G_{x,\Gamma}(N)]
  \end{align*}
  અહીં પણ~$\Gamma$ એ શૂન્યથી સૂચકાંકિત ચલોની યાદી છે,
  અને~$\Gamma[n]$ એ~$n$મા સ્થાને રહેલો ચલ દર્શાવે છે.
  દાખલા તરીકે, $\Gamma$ એ~$x,y,z$ હોય તો~$\Gamma[1]$
  એ~$y$ છે.

  છેલ્લા સમીકરણમાં ચલ~$x$ એવો કોઈપણ ચલ પસંદ કરાય છે જે
  $\Gamma$માં નથી.
\end{defn}

હવે સાબિતી આપ્યા વિના કેટલાંક પરિણામો નોંધીએ:

\begin{prop}
  $M \aconvone M'$ હોય અને~$\Gamma$ એ~$\FV{M}$ને સમાવી
  લેતી કોઈપણ યાદી હોય તો
  $F_\Gamma(M) \eqs F_\Gamma(M')$.
\end{prop}

\end{document}
